\section*{Supplemental Tables}

% Stats for the different stages of the assembly
\begin{table}[ht]
    \centering
    \suptablecaption{\textbf{Assembly statistics for the \balgla{} genome
        at different stages of the assembly process.}
        Default contigs refers to the sequences assembled by \textsc{hifiasm} with
            default parameters.
        Optimized contigs denotes the sequences generated using \textsc{hifiasm}
            followed by removal of contaminant sequences and purging of haplotigs.
        Hi-C scaffolds refer to the scaffolds generated by \textsc{YaHS}.
        Curated assembly refers to the sequences generated after manual curation of 
            the Hi-C contact map and polishing with \textsc{Inspector}.
        \textsc{BUSCO} results refer to the \texttt{arthropoda\_odb10} dataset.}
    \begin{tabular}{lrrrr}
        \toprule
        Assembly Stage &
            \shortstack{Default \\ contigs} &
            \shortstack{Optimized \\ contigs} &
            \shortstack{Hi-C \\ scaffolds} &
            \shortstack{Curated \\ assembly} \\
        \midrule
        Sequence length (Mbp)            & 1.64  & 1.05  & 1.05  & 1.04  \\
        Fragments (N)                    & 4,212 & 1,342 & 650   & 592   \\
        Contig N50 (Mbp)                 & 0.734 & 1.35  & 1.18  & 1.18  \\
        Scaffold N50 (Mbp)               & NA    & NA    & 50.96 & 51.41 \\
        \textsc{BUSCO}$^5$ Complete (\%) & 95.66 & 93.88 & 93.98 & 94.08 \\
        \textsc{BUSCO} Single (\%)       & 45.31 & 89.04 & 89.24 & 89.44 \\
        \textsc{BUSCO} Duplicates (\%)   & 50.35 & 4.84  & 4.74  & 4.64  \\
        \textsc{BUSCO} Fragmented (\%)   & 0.79  & 1.18  & 0.89  & 0.89  \\
        \textsc{BUSCO} Missing (\%)      & 3.55  & 4.94  & 5.13  & 5.03  \\
        \textsc{BUSCO} size (N)          & 1,013 & 1,013 & 1,013 & 1,013 \\
        \bottomrule
    \end{tabular}
    \label{tab:sup_asm_stats}
\end{table}

% Gene-completeness in the annotation
\begin{table}[ht]
    \centering
    \suptablecaption{\textbf{Completeness metrics for the \balgla{} protein-coding
        gene annotation.}
        \textsc{BUSCO} completeness was determined by \textsc{compleasm}
        \texttt{protein}.}
    \begin{tabular}{lrr}
        \toprule
        Metric &
            \texttt{arthropoda\_odb10} &
            \texttt{crustacea\_odb12} \\
        \midrule
        \textsc{BUSCO} Complete (\%)    & 93.58 & 84.77 \\
        \textsc{BUSCO} Single (\%)      & 79.76 & 72.66 \\
        \textsc{BUSCO} Duplicates (\%)  & 13.82 & 12.11 \\
        \textsc{BUSCO} Fragmented (\%)  & 1.78  & 4.10  \\
        \textsc{BUSCO} Missing (\%)     & 4.64  & 11.13 \\
        \textsc{BUSCO} Dataset Size (N) & 1,013 & 1,536 \\
        \bottomrule
    \end{tabular}
    \label{tab:sup_annot_stats}
\end{table}

% Repeat annotation
\begin{table}[ht]
    \centering
    \suptablecaption{\textbf{High-level repeat annotation of the \balgla{}
        assembly.} Repeats were annotated with the \textsc{EarlGrey} software.}
    \begin{tabular}{lrrrr}
        \toprule
        Classification &
            \shortstack{Number of \\ sites (bp)} &
            \shortstack{Count of \\ elements} &
            \shortstack{Proportion of \\ the genome} &
            \shortstack{Number of \\ distinct elements} \\
        \midrule
        DNA                       & 88,203,702  & 217,040 & 0.08 & 335   \\
        LINE                      & 95,553,328  & 199,835 & 0.09 & 511   \\
        LTR                       & 39,476,928  & 87,021  & 0.04 & 362   \\
        \shortstack[l]{Other (Simple Repeat, \\
            ~~Microsatellite, RNA)} & 37,933,124  & 77,043  & 0.04 & 112   \\
        Penelope                  & 25,187,297  & 61,172  & 0.02 & 112   \\
        Rolling Circle            & 6,897,773   & 17,411  & 0.01 & 45    \\
        SINE                      & 828,944     & 2,039   & 0.00 & 6     \\
        Unclassified              & 404,513,462 & 966,499 & 0.39 & 1,520 \\
        \bottomrule
    \end{tabular}
    \label{tab:sup_earlgrey_summary}
\end{table}

% Stats about polymorphism in the reference individual and the popgen samples.
\begin{table}[ht]
    \centering
    \suptablecaption{\textbf{Genetic variation in central Oregon \balgla{}.}}
    \begin{tabular}{lrr}
        \toprule
        Metric &
            Reference Individual &
            Cape Perpetua N=6 \\
        \midrule
        Reference sequence length (bp)      & 904,940,360 & 904,940,360 \\
        Genotyped sites post-filtering (bp) & 314,105,136 & 111,668,206 \\
        Total variant sites (bp)            & 9,081,889   & 21,395,217  \\
        Percent variant sites (\%)          & 2.89        & 19.16       \\
        SNPs (N)                            & 8,725,877   & 20,501,788  \\
        Indels (N)                          & 356,012     & 893,429     \\
        Multiallelic SNPs (N)               & NA          & 996,103     \\
        Percent multiallelic SNPs (\%)      & NA          & 4.86        \\
        Mean variant sites per kbp (N)      & 14.66       & 41.05       \\
        Median variant sites per kbp (N)    & 6           & 29          \\
        St Dev variant sites per kbp (N)    & 19.86       & 40.82       \\
        Variant sites in CDSs (N)           & 398,773     & 1,118,874   \\
        Percent variants in CDS (\%)        & 4.69        & 5.23        \\
        Mean variants per CDS (N)           & 4.34        & 12.19       \\
        Median variants per CDS (N)         & 1           & 6           \\
        St Dev variants per CDS (N)         & 11.75       & 25.84       \\
        \bottomrule
    \end{tabular}
\label{tab:sup_popgen_stats}
\end{table}

% Stats about the phasCons conserved elements and enrichment of genomic elements.
\begin{sidewaystable}[ht]
    \suptablecaption{\textbf{Enchrichment of annotated features among \textsc{phastCons}
        highly conserved elements.}
        The total feature proportion was determined according to the 841,166,098 bp
        corresponding to the 16 chromosome-level scaffolds.
        Proportion of highly-conserved elements (HCE) was calculated based on the
        50,849,898 bp-span of all HCEs in the genome.
        The number of chromosomes denote the number of chromosome-level scaffolds
        containing annotations of the given feature.}
    \centering
    \label{tab:sup_phastcons_enrichment}
    \begin{tabular}{lrrrrrrrr}
    \toprule
        Feature &
            \shortstack{Total feature \\ length (bp)} &
            \shortstack{Total feature \\ proportion} &
            \shortstack{HCE feature \\ length (bp)} &
            \shortstack{HCE feature \\ proportion} &
            \shortstack{Mean $\log_2$ \\ enrichment} &
            \shortstack{Median $\log_2$ \\ enrichment} &
            \shortstack{St Dev $\log_2$ \\ enrichment} &
            \shortstack{Number of \\ chromosomes} \\
    \midrule
        3' UTR     & 2,358,104   & $2.80\mathrm{e}{-3}$ & 152,598    & $3.00\mathrm{e}{-3}$ & $-0.0861$    & 0.0332       & 0.508        & 16 \\
        5' UTR     & 530,151     & $6.30\mathrm{e}{-4}$ & 27,679     & $5.44\mathrm{e}{-4}$ & $-0.292$     & $-0.332$     & 0.453        & 16 \\
        CDS        & 16,273,411  & 0.0193               & 2,225,138  & 0.0438               & 0.752        & 0.984        & 0.926        & 16 \\
        TE         & 370,573,528 & 0.441                & 17,505,920 & 0.344                & $-0.359$     & $-0.305$     & 0.180        & 16 \\
        intergenic & 196,612,776 & 0.234                & 10,688,761 & 0.210                & $-0.167$     & $-0.178$     & 0.167        & 16 \\
        intron     & 111,562,654 & 0.133                & 8,575,960  & 0.169                & 0.419        & 0.297        & 0.432        & 16 \\
        lncRNA     & 2,211,027   & $2.63\mathrm{e}{-3}$ & 130,744    & $2.57\mathrm{e}{-3}$ & $-0.174$     & $-0.133$     & 0.567        & 16 \\
        multiple   & 141,022,121 & 0.168                & 11,537,936 & 0.227                & 0.418        & 0.438        & 0.167        & 16 \\
        rRNA       & 2,106       & $2.50\mathrm{e}{-6}$ & 0.000      & 0.000                & \texttt{NaN} & \texttt{NaN} & \texttt{NaN} & 6  \\
        tRNA       & 20,220      & $2.40\mathrm{e}{-5}$ & 5,162      & $1.02\mathrm{e}{-4}$ & 2.143        & 2.041        & 0.959        & 16 \\
   \bottomrule
\end{tabular}
\end{sidewaystable}

% Reference for the allozyme genes in our assembly.
\begin{sidewaystable}[ht]
    \suptablecaption{\textbf{Classic allozyme loci identified in the \balgla{}
        genome assembly.}
        Homology between the reference allozyme loci and the \bgla{} transcripts
        was determined according to the best \textsc{BLAST} hit.}
    \centering
    \begin{tabular}{lllrrllll}
        \toprule
        \shortstack{\bgla{} \\ transcript ID} &
            Locus &
            Chromosome &
            Start (bp) &
            End (bp) & 
            Gene description &
            \shortstack{NCBI \\ accession} &
            Reference \\
        \midrule
        \texttt{mrna-5520} & \textit{Adh} & chr04 & 19{,}904{,}033 & 19{,}930{,}487 & 
            alcohol dehydrogenase &
            \texttt{QIA50413.1} &~\cite{nunez_footprints_2020} \\
        \texttt{mrna-15301} & \textit{Gpi} & chr12 & 15{,}750{,}365 & 15{,}777{,}818 &
            \shortstack[l]{glucose phosphate \\ isomerase} &
            \texttt{QIA50414.1} &~\cite{nunez_footprints_2020} \\
        \texttt{mrna-1965} & \textit{Me.1} & chr02 & 5{,}827{,}144 & 5{,}858{,}664 &
            \shortstack[l]{NADP-dependent \\ malic enzyme} &
            \texttt{KAF0299259.1} &~\cite{kim_draft_2019} \\
        \texttt{mrna-3861} & \textit{Me.2} & chr03 & 12{,}307{,}597 & 12{,}337{,}955 &
            \shortstack[l]{NADP-dependent \\ malic enzyme} &
            \texttt{KAF0308712.1} &~\cite{kim_draft_2019} \\
        \texttt{mrna-5994} & \textit{Mpi} & chr04 & 40{,}166{,}344 & 40{,}173{,}580 &
            \shortstack[l]{mannose-6-phosphate \\ isomerase} &
            \texttt{QIA50417.1} &~\cite{nunez_footprints_2020} \\
        \texttt{mrna-5969} & \textit{Pgm} & chr04 & 38{,}875{,}727 & 38{,}886{,}668 &
            phosphoglucomutase &
            \texttt{QIA50416.1} &~\cite{nunez_footprints_2020} \\
        \texttt{mrna-17464} & \textit{Sod} & chr14 & 42{,}896{,}065 & 42{,}935{,}242 &
            superoxide dismutase &
            \texttt{QIA50415.1} &~\cite{nunez_footprints_2020} \\

        \bottomrule
    \end{tabular}
    \label{tab:sup_allozymes}
\end{sidewaystable}

% dN/dS for the allozymes.
\begin{sidewaystable}[ht]
    \suptablecaption{\textbf{Patterns of \dnds{} in the \balgla{} classic allozyme loci.}
        Substitutions were calculated based on the codon alignment between \bgla{}
            and \textit{B. crenatus} orthologs.
        All loci showed \dnds{} $<1$, indicating evolution under a regime of
            purifying selection.
        N\textsubscript{C} = number of non-synonymous sites;
        S\textsubscript{C} = number of synonymous sites;
        N\textsubscript{D} = number of non-synonymous differences;
        S\textsubscript{D} = number of synonymous differences;
        p\textsubscript{N} = proportion of non-synonymous differences;
        p\textsubscript{S} = proportion of synonymous differences;
        d\textsubscript{N} = Jukes-Cantor corrected proportion of
            non-synonymous differences;
        d\textsubscript{S} = Jukes-Cantor corrected proportion of
            synonymous differences;
    }
    \centering
    \begin{tabular}{llrrrrrrrrrr}
        \toprule
        \shortstack{\bgla{} \\ transcript ID} &
            Locus &
            \shortstack{Transcript \\ length (bp)} &
            N\textsubscript{C} & % Nc = NonSyn Sites
            S\textsubscript{C} & % Sc = Syn Sites
            N\textsubscript{D} & % Nd = NonSyn Differences
            S\textsubscript{D} & % Sd = Syn Differences
            p\textsubscript{N} & % pN = Proportion NonSyn
            p\textsubscript{S} & % pS = Proportion Syn
            d\textsubscript{N} & % dN = Adj Proportion NonSyn
            d\textsubscript{S} & % dS = Adj Proportion Syn
            \dnds{} \\           % dN/dS
        \midrule
            % TransID             Locus          length      Nc            Sc        Nd       Sd        pN       pS      dN       dS   dN/dS
            \texttt{mrna-5520}  & \textit{Adh}  & 1{,}140 & 858.167     & 272.833 & 15     & 64      & 0.0175 & 0.235 & 0.0177 & 0.281 & 0.0629 \\
            \texttt{mrna-15301} & \textit{Gpi}  & 1{,}758 & 1{,}310.667 & 399.333 & 47.833 & 101.167 & 0.0365 & 0.253 & 0.0374 & 0.309 & 0.121  \\
            \texttt{mrna-5994}  & \textit{Mpi}  & 1{,}233 & 907.167     & 322.833 & 39     & 79      & 0.0430 & 0.245 & 0.0443 & 0.296 & 0.149  \\
            \texttt{mrna-1965}  & \textit{Me.1} & 1{,}818 & 1{,}376.667 & 438.333 & 37.500 & 113.500 & 0.0272 & 0.259 & 0.0277 & 0.318 & 0.0874 \\
            \texttt{mrna-3861}  & \textit{Me.2} & 1{,}902 & 1{,}439.167 & 459.833 & 48     & 130     & 0.0334 & 0.283 & 0.0341 & 0.355 & 0.0961 \\
            \texttt{mrna-5969}  & \textit{Pgm}  & 1{,}815 & 1{,}287.833 & 404.167 & 32.167 & 93.833  & 0.0250 & 0.232 & 0.0254 & 0.279 & 0.0914 \\
            \texttt{mrna-17464} & \textit{Sod}  &     605 & 473.167     & 150.833 & 27     & 43      & 0.0571 & 0.285 & 0.0593 & 0.359 & 0.165  \\
        \bottomrule
    \end{tabular}
    \label{tab:sup_allozymes_dnds}
\end{sidewaystable}

% MKT for the allozymes
\begin{sidewaystable}[ht]
    \suptablecaption{\textbf{McDonald-Kreitman test on the \balgla{} classic allozyme
        loci.}
        MK test calculated by \textsc{MKado} based on the codon-alignment between
            orthologous \textit{B. crenatus} outgroup sequences and polymorphic 
            \bgla{} ingroup sequences.
        Both \textit{Gpi} and \textit{Me.2} showed a significant difference in the
            number vs expected synonymous and non-synonymous divergence and polymorphism.
        D\textsubscript{N} = non-synonymous substitutions;
        D\textsubscript{S} = synonymous substitutions;
        P\textsubscript{N} = non-synonymous polymorphisms;
        P\textsubscript{S} = synonymous polymorphism;
        $\alpha$ = proportion of non-synonymous substitutions fixed by
            positive selection;
        NI = Neutrality Index;
        DoS = Direction of Selection.
    }
    \centering
    \begin{tabular}{llrrrrrrrrr}
        \toprule
        \shortstack{\bgla{} \\ transcript ID} &
            Locus &
            \shortstack{Transcript \\ length (bp)} &
            D\textsubscript{N} &        % Dn, Fixed NonSyn differences
            D\textsubscript{S} &        % Ds, Fixed Syn differences
            P\textsubscript{N} &        % Pn, Polymorphic NonSyn sites
            P\textsubscript{S} &        % Ps, Polymorphic Syn sites
            \shortstack{Corrected \\ \textit{p}-value} &
            $\alpha$ &
            \ NI &
            \ DoS \\
        \midrule
            % Bg trans            Gene            length    Dn   Ds    Pn   Ps    Pval                 alpha                    NI      DoS
            \texttt{mrna-5520}  & \textit{Adh}  & 1{,}140 & 13 & 46  & 7  & 32  & 0.932              & 0.226                  & 0.774 & 0.041                  \\
            \texttt{mrna-15301} & \textit{Gpi}  & 1{,}758 & 36 & 76  & 7  & 62  & \textbf{2.112e--3} & 0.762                  & 0.239 & 0.220                  \\
            \texttt{mrna-1965}  & \textit{Me.1} & 1{,}818 & 27 & 105 & 8  & 31  & 1.00               & $-3.584\mathrm{e}{-3}$ & 1.004 & $-5.830\mathrm{e}{-4}$ \\
            \texttt{mrna-3861}  & \textit{Me.2} & 1{,}902 & 37 & 81  & 15 & 117 & \textbf{1.022e--3} & 0.719                  & 0.281 & 0.200                  \\
            \texttt{mrna-5994}  & \textit{Mpi}  & 1{,}233 & 19 & 63  & 19 & 40  & 0.443              & $-0.575$               & 1.575 & $-0.090$               \\
            \texttt{mrna-5969}  & \textit{Pgm}  & 1{,}815 & 23 & 73  & 10 & 65  & 0.274              & 0.512                  & 0.488 & 0.106                  \\
            \texttt{mrna-17464} & \textit{Sod}  &     605 & 22 & 39  & 6  & 13  & 0.932              & 0.182                  & 0.818 & 0.045                  \\
        \bottomrule
    \end{tabular}
    \label{tab:sup_allozymes_mkt}
\end{sidewaystable}

% HKA test per-locus vs pooled remainder
\begin{table}[ht]
    \centering
    \suptablecaption{\textbf{Classic Hudson-Kreitman-Aguad\'{e} test for the seven
        classic allozyme loci, testing each locus against the pooled remainder.}
        Polymorphism was measured as segregating silent sites, pooling synonymous
        sites within CDS with all alignable sites in introns and 5\,kbp flanks per
        the 1991 convention. Divergence was measured against the \textit{B.
        crenatus} ortholog across the same sites. $T+1$ was jointly estimated as
        8.15 across all seven loci. Only \textit{Me.1} reached significance, with
        a deficit of polymorphism relative to divergence.}
    \begin{tabular}{l l r r r r r r l}
        \toprule
        \shortstack{\bgla{} \\ transcript ID} &
            Locus &
            obs\textsubscript{poly} &
            exp\textsubscript{poly} &
            obs\textsubscript{div}  &
            exp\textsubscript{div}  &
            $\chi^2$ &
            \textit{p}-value &
            Direction \\
        \midrule
        \texttt{mrna-5520}  & \textit{Adh}  &  439    &  544    & 1{,}548 & 1{,}458 & 0.54 & 0.461 & deficit \\
        \texttt{mrna-15301} & \textit{Gpi}  & 1{,}651 & 1{,}403 & 3{,}461 & 3{,}764 & 0.52 & 0.471 & excess  \\
        \texttt{mrna-1965}  & \textit{Me.1} &  441    &  989    & 3{,}198 & 2{,}650 & 4.59 & \textbf{0.032} & deficit \\
        \texttt{mrna-3861}  & \textit{Me.2} & 2{,}017 & 1{,}586 & 3{,}821 & 4{,}252 & 1.15 & 0.284 & excess  \\
        \texttt{mrna-5994}  & \textit{Mpi}  &  712    &  861    & 2{,}457 & 2{,}308 & 0.45 & 0.505 & deficit \\
        \texttt{mrna-5969}  & \textit{Pgm}  &  861    &  776    & 1{,}995 & 2{,}080 & 0.18 & 0.673 & excess  \\
        \texttt{mrna-17464} & \textit{Sod}  & 2{,}553 & 2{,}526 & 6{,}744 & 6{,}771 & 0.00 & 0.964 & excess  \\
        \bottomrule
    \end{tabular}
    \label{tab:sup_hka}
\end{table}

% Breakdown of the per-feature Pi
\begin{sidewaystable}[ht]
    \centering
    \suptablecaption{\textbf{Nucleotide diversity ($\pi$) across different
        genomic features}}
    \label{tab:sup_feature_pi}
    \begin{tabular}{lrrrrrrrrrrr}
        \toprule
        Feature &
            \shortstack{Mean \\ $\pi$} &
            \shortstack{Median \\ $\pi$} &
            \shortstack{StDev \\ $\pi$} &
            \shortstack{Number \\ of \\ windows} &
            \shortstack{Mean \\ window \\ length \\ (bp)} &
            \shortstack{Median \\ window \\ length \\ (bp)} &
            \shortstack{StDev \\ window \\ length \\ (bp)} &
            \shortstack{Number \\ of \\ sites} &
            \shortstack{Mean \\ sites per \\ window} &
            \shortstack{Median \\ sites per \\ window} &
            \shortstack{StDev \\ sites per \\ window} \\
        \midrule
        % Feature                            meanPi   medianPi               sdPi    nWinds  meanWinLen  MedWinLen sdWinLen     nBP          meanWinBP  medWinBP  sdWinBP
        \shortstack[l]{Genome\\wide}       & 0.0506 & 0.0531               & 0.0188 & 44,599 & 9,998.570 & 10,000 & 103.774   & 93,583,599 & 2,098.334 & 1,932  & 821.464   \\
        Intergenic                         & 0.0507 & 0.0526               & 0.0210 & 29,219 & 7,688.054 & 10,000 & 3531.991  & 45,677,153 & 1,563.269 & 1,431  & 956.166   \\
        CDS                                & 0.0302 & 0.0270               & 0.0221 & 78,630 & 270.688   & 161    & 405.314   & 9,447,233  & 120.148   & 80     & 163.758   \\
        Intron                             & 0.0646 & 0.0650               & 0.0299 & 67,638 & 2,182.716 & 686    & 5,257.832 & 34,345,262 & 507.781   & 209    & 1,090.025 \\
        5' UTR                             & 0.0398 & 0.0338               & 0.0309 & 6,013  & 153.063   & 114    & 148.787   & 424,412    & 70.582    & 54     & 63.678    \\
        3' UTR                             & 0.0489 & 0.0468               & 0.0294 & 8,218  & 648.360   & 542    & 485.406   & 2,083,933  & 253.582   & 210    & 194.076   \\
        lncRNA                             & 0.0341 & 0.0302               & 0.0246 & 27,168 & 399.958   & 184    & 511.811   & 4,557,588  & 167.756   & 91     & 200.540   \\
        tRNA                               & 0.0118 & $5.21\mathrm{e}{-3}$ & 0.0236 & 430    & 66.177    & 72     & 16.982    & 13,804     & 32.102    & 32     & 14.305    \\
        \shortstack[l]{0-fold\\degenerate} & 0.0107 & $2.92\mathrm{e}{-3}$ & 0.0178 & 75,895 & 269.321   & 162    & 401.035   & 5,397,246  & 71.115    & 46     & 96.221    \\
        \shortstack[l]{4-fold\\degenerate} & 0.0790 & 0.0750               & 0.0483 & 39,170 & 281.364   & 176    & 380.879   & 1,410,361  & 36.006    & 23     & 46.585    \\
    \bottomrule
    \end{tabular}
\end{sidewaystable}

% Genebank/RefSeq data used
\begin{table}[ht]
    \centering
    \suptablecaption{\textbf{Publicly available NCBI genomes used for comparative
        analyses}}
    \begin{tabular}{llll}
        \toprule
            Species & NCBI accession & Reference &  Note \\
        \midrule
            \textit{Amphibalanus amphitrite} & \texttt{GCF\_019059575.1} & \cite{kim_draft_2019}               & \shortstack[l]{RefSeq genome\\and annotation} \\
            \textit{Amphibalanus amphitrite} & \texttt{GCA\_037642225.1} & \cite{han_new_2024}                 & \shortstack[l]{GenBank genome\\without annotation} \\
            \textit{Capitulum mitella}       & \texttt{GCA\_030062745.1} & \cite{chen_capmit_2021}             & \shortstack[l]{GenBank genome\\without annotation} \\
            \textit{Pollicipes pollicipes}   & \texttt{GCF\_011947565.3} & \cite{bernot_chromosome-level_2022} & \shortstack[l]{RefSeq genome\\and annotation} \\
            \textit{Drosophila melanogaster} & \texttt{GCF\_000001215.4} & \cite{adams_drosophila_2000}        & \shortstack[l]{RefSeq genome\\and annotation} \\
        \bottomrule
    \end{tabular}
    \label{tab:sup_ncbi_genomes}
\end{table}

% SRA data
\begin{table}[ht]
    \centering
    \suptablecaption{\textbf{Publicly available NCBI short-read archive data used
        for comparative analyses}}
    \begin{tabular}{lll}
        \toprule
            SRA & BioProject & Species \\
        \midrule
            % SRA                  BioProject             Species
            \texttt{SRR8775109}  & \texttt{PRJNA528777} & \textit{Amphibalanus improvisus} \\
            \texttt{SRR8775110}  & \texttt{PRJNA528777} & \textit{Amphibalanus improvisus} \\
            \texttt{SRR8775111}  & \texttt{PRJNA528777} & \textit{Amphibalanus improvisus} \\
            \texttt{SRR8775112}  & \texttt{PRJNA528777} & \textit{Amphibalanus improvisus} \\
            \texttt{SRR8775113}  & \texttt{PRJNA528777} & \textit{Amphibalanus improvisus} \\
            \texttt{SRR8775114}  & \texttt{PRJNA528777} & \textit{Amphibalanus improvisus} \\
            \texttt{SRR8775115}  & \texttt{PRJNA528777} & \textit{Amphibalanus improvisus} \\
            \texttt{SRR8775116}  & \texttt{PRJNA528777} & \textit{Amphibalanus improvisus} \\
            \texttt{SRR8775117}  & \texttt{PRJNA528777} & \textit{Amphibalanus improvisus} \\
            \texttt{SRR8775118}  & \texttt{PRJNA528777} & \textit{Amphibalanus improvisus} \\
            \texttt{SRR8775119}  & \texttt{PRJNA528777} & \textit{Amphibalanus improvisus} \\
            \texttt{SRR8775120}  & \texttt{PRJNA528777} & \textit{Amphibalanus improvisus} \\
            \texttt{SRR8775121}  & \texttt{PRJNA528777} & \textit{Amphibalanus improvisus} \\
            \texttt{SRR8775122}  & \texttt{PRJNA528777} & \textit{Amphibalanus improvisus} \\
            \texttt{SRR8775123}  & \texttt{PRJNA528777} & \textit{Amphibalanus improvisus} \\
            \texttt{SRR8775124}  & \texttt{PRJNA528777} & \textit{Amphibalanus improvisus} \\
            \texttt{SRR8775125}  & \texttt{PRJNA528777} & \textit{Amphibalanus improvisus} \\
            \texttt{SRR8775126}  & \texttt{PRJNA528777} & \textit{Amphibalanus improvisus} \\
            \texttt{SRR8775127}  & \texttt{PRJNA528777} & \textit{Amphibalanus improvisus} \\
            \texttt{SRR8775128}  & \texttt{PRJNA528777} & \textit{Amphibalanus improvisus} \\
            \texttt{SRR8775129}  & \texttt{PRJNA528777} & \textit{Amphibalanus improvisus} \\
            \texttt{SRR8775130}  & \texttt{PRJNA528777} & \textit{Amphibalanus improvisus} \\
            \texttt{SRR8775131}  & \texttt{PRJNA528777} & \textit{Amphibalanus improvisus} \\
            \texttt{SRR8755479}  & \texttt{PRJNA528169} & \textit{Amphibalanus improvisus} \\
            \texttt{SRR18959849} & \texttt{PRJNA816681} & \textit{Capitulum mitella} \\
            \texttt{SRR14354747} & \texttt{PRJNA725059} & \textit{Capitulum mitella} \\
            \texttt{SRR20067441} & \texttt{PRJNA856037} & \textit{Capitulum mitella} \\
            \texttt{SRR20067442} & \texttt{PRJNA856037} & \textit{Capitulum mitella} \\
            \texttt{SRR20067443} & \texttt{PRJNA856037} & \textit{Capitulum mitella} \\
            \texttt{SRR20067446} & \texttt{PRJNA856037} & \textit{Capitulum mitella} \\
            \texttt{SRR20067447} & \texttt{PRJNA856037} & \textit{Capitulum mitella} \\
            \texttt{SRR20067448} & \texttt{PRJNA856037} & \textit{Capitulum mitella} \\
            \texttt{SRR20067449} & \texttt{PRJNA856037} & \textit{Capitulum mitella} \\
            \texttt{SRR20067444} & \texttt{PRJNA856037} & \textit{Capitulum mitella} \\
            \texttt{SRR20067445} & \texttt{PRJNA856037} & \textit{Capitulum mitella} \\
            \texttt{SRR13527601} & \texttt{PRJNA681321} & \textit{Capitulum mitella} \\
            \texttt{SRR13527602} & \texttt{PRJNA681321} & \textit{Capitulum mitella} \\
            \texttt{SRR13527603} & \texttt{PRJNA681321} & \textit{Capitulum mitella} \\
            \texttt{SRR13527604} & \texttt{PRJNA681321} & \textit{Capitulum mitella} \\
        \bottomrule
    \end{tabular}
    \label{tab:sup_sra_data}
\end{table}

